\section{Calculation of the Wilson coefficients}\label{sec:calc}

\subsection{Method of Projectors}

To compute the coefficients $\zeta_{ij}(t)$ we use the so-called
``method of projectors''\,\cite{Gorishnii:1983su,Gorishnii:1986gn},
which consists of constructing external states $|k\rangle$ and differential
operators $D_k$ for which
\begin{equation}
  \begin{split}
    P_{k}[\calo_{i} (x) ] \equiv D_{k}\langle
    0|\calo_{i} (x) |k\rangle =
    \delta_{ik}\,,
    \label{eq:projectors}
  \end{split}
\end{equation}
where we have dropped the Lorentz indices for convenience, and we define
the matrix element to include only diagrams which are one-particle
irreducible (1\abbrev{PI}) with respect to (w.r.t.)
\qcd\ particles. Applying $P_{k}$ on both sides of \eqn{eq:zetas}, one
obtains
\begin{equation}
  \begin{split}
    P_{k}[\tcalo_{i}(t,x)] =
    \zeta_{ij}(t)P_{k}[\calo_{j}(x)]\,.
    \label{eq:project}
  \end{split}
\end{equation}
Since the $\zeta_{ij}(t)$ only depend on the flow time $t$ and the
renormalization scale $\mu$, we can choose arbitrary values for all
other dimensional parameters in this equation. Setting them to zero
turns all higher-order corrections on the r.h.s.\ into massless
tadpoles, so that \eqn{eq:projectors} is only required to hold at
tree-level.  One thus obtains
\begin{equation}
  \begin{split}
    \zeta_{ij}(t) = P_j[\tcalo_i(t,x)]\Big|_{p=m=0}\,,
  \end{split}
\end{equation}
where $m$ and $p$ collectively denote all masses and external momenta.
The right-hand side thus results in vacuum diagrams whose only dimensional
scale is $t$.

In order to find suitable projectors, we first derive the Feynman rules
for the relevant terms of the operators. For example,\footnote{ All
  Feynman diagrams in this paper were drawn using Ti\textit{k}Z-Feynman
  \cite{Ellis:2016jkw}.}
\begin{align}
  \calo_{1,\mu\nu} &= \partial_\mu A_\rho^c\partial_\nu A_\rho^c +
  \cdots&\Rightarrow&&
  \begin{gathered}
	\includegraphics{images/O1.pdf}
  \end{gathered} =
  - p_{1,\mu} p_{2,\nu} \delta_{\alpha\beta}\delta^{ab},
\end{align}
where the momenta are defined to be outgoing.
This suggests to use
\begin{equation}
  \begin{split}
    P_{1}[X] = - \frac{\delta^{ab}}{N_A} P_{\alpha \beta |
      \rho \mu | \sigma \nu}
    \deriv{}{p_{1,\rho}}
    \deriv{}{p_{2,\sigma}}
    \langle 0 | A_\alpha^a (p_1) A_\beta^b (p_2) X_{\mu\nu}| 0 \rangle\,,
    \label{eq:p1}
  \end{split}
\end{equation}
where $N_A$ is the dimension of the adjoint representation of the gauge
group; for SU($N_c$), it is $N_A=N_c^2-1$. The projector onto the Lorentz structure is defined by
\begin{equation}
  \begin{split}
    P_{\alpha_1\beta_1|\cdots|\alpha_n\beta_n}
    T_{\alpha_1\beta_1\cdots\alpha_n\beta_n} = \left\{\begin{array}{ll}
    1\phantom{\Bigg(}\quad &\text{for}\quad
    T_{\alpha_1\beta_1\cdots\alpha_n\beta_n}=\delta_{\alpha_1\beta_1}
    \cdots\delta_{\alpha_n\beta_n}\,,\\[.3em]
    0& \text{for any other}\\[-.1em]
    & \text{linearly independent}\\[-.1em]
    & \text{Lorentz tensor.}
    \end{array}\right.
  \end{split}
\end{equation}

In the appendix, one can find the relevant parts of the Feynman rules
for the other operators, which in a similar way lead to the projectors
\begin{equation}
  \begin{split}
    P_2[X] &=  - \frac{\delta^{ab}}{4 N_A} P_{\alpha \beta
      | \mu \nu | \rho \sigma}
    \deriv{}{p_{1,\rho}}
    \deriv{}{p_{2,\sigma}}
    \langle 0 | A_\alpha^a (p_1) A_\beta^b (p_2) X_{\mu \nu} | 0
    \rangle \,,\\
    P_{3_f}[X] &= -i \frac{\delta^{i j}}{4 N_c} P_{\rho
      \mu | \sigma \nu} \deriv{}{p_{2,\sigma}} \text{Tr} \left[
      \gamma_{\rho}
    \langle 0 | \psi_f^j (p_2) \bar{\psi}_f^i (p_1) X_{\mu \nu}
    | 0 \rangle \right]\,,\\
    P_{4_f}[X] &= -i \frac{\delta^{i j}}{4 N_c} P_{\mu \nu
      | \sigma \rho} \deriv{}{p_{2,\sigma}} \text{Tr} \left[
      \gamma_{\rho} \langle 0 | \psi_f^j (p_2) \bar{\psi}_f^i (p_1)
      X_{\mu \nu} | 0 \rangle \right] \\
    & \qquad - \frac{1}{2} \frac{\delta^{i j}}{4 N_c} P_{\mu
      \nu} \deriv{}{m_0} \text{Tr} \left[ \langle 0 | \psi_f^j (p_2)
      \bar{\psi}_f^i (p_1) X_{\mu \nu} | 0 \rangle \right]\,,
    \label{eq:pj}
  \end{split}
\end{equation}
where $N_c$ is the dimension of the fundamental representation of the
gauge group, i.e.\ the number of colors, and $i$ and $j$ are the corresponding indices.
The trace appearing in the projectors $P_{3_f}$ and
$P_{4_f}$ is taken w.r.t.\ the spinor indices of the Green's function,
and $f$ denotes the associated quark flavor.
Note that $P_{4_f}$ is constructed such that
\begin{equation}
  \begin{split}
    P_{4_f}[\calo_4 + 2\calo_5] = 0
  \end{split}
\end{equation}
in order to ensure that only those fermionic operators are taken into
account which do not vanish according to the \eom, see \eqn{eq:eom}.

With this procedure, we get a mixing matrix which distinguishes between
different quark flavors. To avoid confusion with $\zeta_{ij} (t)$, which
is the mixing matrix between operators summed over all flavors, we will
call it $\Omega_{ij} (t)$. This matrix is then defined by
\begin{equation}
	\tcalo_{i, \mu \nu} (t, x) = \Omega_{i j} (t) \calo_{j, \mu \nu} (x) \, ,
\end{equation}
where double indices in this expression are summed over $\{1, 2,
3_1,\ldots, 3_{n_F}, 4_1,\ldots,4_{n_F}\}$ (see \eqn{eq:ops}), in
contrast to \eqn{eq:emtflow}, where double indices are summed over $\{1,
2, 3, 4\}$. Its general structure is given by
\begin{equation}
	\Omega =
	\begin{pmatrix}
		\Omega_{1 1} & \Omega_{1 2} & \underline{\Omega_{1 3}}^T & \underline{\Omega_{1 4}}^T \\
		\Omega_{2 1} & \Omega_{2 2} & \underline{\Omega_{2 3}}^T & \underline{\Omega_{2 4}}^T \\
		\underline{\Omega_{3 1}} & \underline{\Omega_{3 2}} & \underline{\underline{\Omega_{3 3}}} & \underline{\underline{\Omega_{3 4}}} \\
		\underline{\Omega_{4 1}} & \underline{\Omega_{4 2}} & \underline{\underline{\Omega_{4 3}}} & \underline{\underline{\Omega_{4 4}}}
	\end{pmatrix} \, ,
	\label{eq:Omega}
\end{equation}
where an element $\Omega_{i j}$ represents the mixing between $\calo_i$
and $\tcalo_j$, taking into account individual flavors. Therefore, an
underlined and a double underlined element denotes an $n_F$-dimensional
vector and an $n_F \times n_F$ dimensional matrix, respectively.  Their
elements describe the mixing between different flavors.  As the quarks
are indistinguishable, the $n_F$ and $n_F^2$ dimensional objects
appearing in \eqn{eq:Omega} can each be described by two independent
parameters, named $\omega_{ij}$ and $\bar{\omega}_{ij}$:
\begin{equation*}
	\begin{split}
		\Omega_{i j} = \omega_{i j} \quad \text{for} \quad i,j <
                3, \qquad
		\underline{\Omega_{i j}} = \omega_{i j}
		\begin{pmatrix}
			1 \\
			\vdots \\
			1
		\end{pmatrix} \quad \text{for} \quad i < 3,\
                j > 2 \quad \text{or} \quad i > 2,\ j < 3\,,
	\end{split}
\end{equation*}
\begin{equation}
	\begin{split}
		\underline{\underline{\Omega_{i j}}} =
		\begin{pmatrix}
			\omega_{i j} & \overline{\omega}_{i j} &
                        \overline{\omega}_{i j} & \dots &
                        \overline{\omega}_{i j} \\
			\overline{\omega}_{i j} & \omega_{i j} &
                        \overline{\omega}_{i j} & \dots &
                        \overline{\omega}_{i j} \\
			\overline{\omega}_{i j} & \overline{\omega}_{i j}
                        & \omega_{i j} & \dots & \overline{\omega}_{i j}
                        \\
			\vdots & \vdots & \vdots & \ddots & \vdots \\
			\overline{\omega}_{i j} & \overline{\omega}_{i
                          j} & \overline{\omega}_{i j} & \dots &
                        \omega_{i j}
		\end{pmatrix} \quad \text{for} \quad i,j > 2\, .
	\end{split}
\end{equation}
Summing over the different flavors occurring in \eqn{eq:Omega}, the
relation between $\Omega (t)$ and $\zeta (t)$ can be easily established
by
\begin{align}
	\zeta_{i j} &= \omega_{i j} & \text{for} &\quad i < 3 \, , \\
	\zeta_{i j} &= n_F \, \omega_{i j} & \text{for} &\quad i > 2, \, j < 3 \, , \\
	\zeta_{i j} &= \omega_{i j} + (n_F - 1) \, \overline{\omega}_{i j} & \text{for} &\quad i > 2, \, j > 2 \, .
\end{align}

\subsection{Computational Methods}

The gradient-flow formalism in perturbation theory can be formulated in
terms of a Lagrangian field theory, where the flow equations
(\ref{eq:flow}) are implemented with the help of Lagrange-multiplier
fields \cite{Luscher:2011bx}. The crucial difference between the regular
\qcd\ Feynman rules and those in the gradient-flow formalism is the
occurrence of exponential factors $\exp(-sp^2)$, where $s$ is a
``flow-time variable'', and $p$ the linear combination of
$D$-dimensional external and/or loop momenta. Vertices involving flowed
fields induce an integration over all positive values of the
corresponding flow-time variable, which is, however, bounded from above
by ``propagators'' of the Lagrange-multiplier fields, since they
introduce step functions of the flow-time variables.

We have implemented the Feynman rules into the program
\texttt{qgraf}\,\cite{Nogueira:1991ex,Nogueira:2006pq}, which generates
the Feynman diagrams for the desired matrix elements. Its output is then
transformed to \texttt{FORM}\,\cite{Vermaseren:2000nd,Kuipers:2012rf}
notation by
\texttt{q2e/exp}\,\cite{Harlander:1997zb,Seidensticker:1999bb}. An
in-house set of \texttt{FORM} routines inserts the Feynman rules,
performs the projections onto the relevant color and Lorentz structures
according to the $P_j$ of \eqn{eq:p1} and \noeqn{eq:pj}, and evaluates
the Dirac and color traces using the \texttt{color}
package\,\cite{vanRitbergen:1998pn}.  The result is then expressed in
terms of a linear combination of integrals whose general form is as
follows:
\begin{equation}
  \begin{split}
    I_l( &(d_1, \dots, d_f ), (b_1, \dots, b_n ), ( a_1, \dots, a_n ))
    \\ &\equiv \frac{1}{\pi^{lD/2}} \, t^{lD/2-\sum_{j=1}^n a_j} \left[\prod_{i=0}^f
     \int_0^1 \mathrm{d}u_i u_i^{d_i} \right]
    \left[\prod_{r=1}^l\int\dd^Dk_r\right]\frac{\exp(-t \sum_{j=1}^n b_j
      q_j^2)}{(q_1^2)^{a_1}\dots (q_n^2)^{a_n}}\,,
    \label{eq:fl}
  \end{split}
\end{equation}
where the $a_i$ and $d_i$ are integers ($d_i\geq 0$), $f$ and $l$ is the
number of flow-time and loop integrations, respectively, the $b_j$ are
polynomials in (rescaled) flow-time parameters $u_i$ and the $q_i$ are
linear combinations of the loop momenta $k_j$. For the problem and the
perturbative order under consideration, it is $0\leq f\leq 4$, $1 \leq
l\leq 2$, and $0 \leq n \leq 3$, respectively.  Note that the projectors
defined in \eqs{eq:p1} and \noeqn{eq:pj} eliminate all dependence on
external momenta and masses, so that, after making a suitable ansatz for
the index structure of the integrals, we only have to evaluate scalar
vacuum integrals.  Using the identities\,\cite{Chetyrkin:1981qh}
\begin{equation}
  \begin{split}
    &\int\dd^D k \left(\deriv{}{k}\cdot q\right) f(k,q,\ldots) =
    0\,,
    \label{eq:ibp}
  \end{split}
\end{equation}
and similar ones for the flow-time integrations,
\begin{equation}
  \begin{split}
    \qquad\int_0^1\dd s \deriv{}{s} f(s,\ldots) =
    f(1,\ldots)-f(0,\ldots)\,,
    \label{eq:ibpflow}
  \end{split}
\end{equation}
one can derive relations among these integrals by explicitly performing
the derivatives in the integrand on the l.h.s. These so-called
``integration-by-parts (\abbrev{IBP}) relations'' were fed to
\texttt{Kira}\,\cite{Maierhoefer:2017hyi} which allowed us to reduce all
occurring integrals to a single master integral at \one-loop level, and
six master integrals at \two-loop level using the Laporta
algorithm\,\cite{Laporta:2001dd}.\footnote{The reduction with
  \texttt{Kira} 1.0 takes about 20 minutes on 8 CPU threads and requires
  less than 13\,GB of RAM.}  Their analytical evaluation is possible
along the lines of \citere{Luscher:2010iy}:
\begin{equation}
  \begin{split}
    I_1((),(2),(0)) &= 2^{-D/2} \,,\phantom{\Bigg[} \\[10pt]
    I_2((),(0,2,2),(0,0,0)) &= 2^{-D} \,,\phantom{\Bigg[} \\[10pt]
    I_2((),(1,1,1),(0,0,0)) &= 3^{-D/2} \,,\phantom{\Bigg[} \\[10pt]
    I_2((),(1,1,1),(1,1,0)) &= \frac{1}{D-2} \Bigg[ - 2 \pi \csc \left( \frac{D \pi}{2} \right) \\
    & \qquad \qquad+ \frac{3^{2 - D/2}}{D - 4} \, {_2F}_1 \left(1, 1; 3 - \frac{D}{2}; \frac{3}{4} \right) \Bigg] \,,\\[10pt]
    I_2((),(0,0,2),(1,1,0)) &= - \frac{2^{3 - D} \pi}{D - 2} \csc \left( \frac{D \pi}{2} \right) \,, \phantom{\Bigg[} \\[10pt]
    I_2((0),(2-u_1,u_1,u_1),(0,0,0)) &= 2^{2 - 2D} B_{1/4} \left( 1 - \frac{D}{2}, 1 - \frac{D}{2} \right) \,,\phantom{\Bigg[} \\[10pt]
    I_2((0),(1+u_1,1+u_1,1-u_1),(0,0,0)) &= 2^{2 - 2 D} \Bigg[ B_{3/4} \left( 1 - \frac{D}{2}, 1 - \frac{D}{2} \right) \\
    & \qquad \qquad - B_{1/2} \left( 1 - \frac{D}{2}, 1 - \frac{D}{2} \right) \Bigg] \,.
    \label{eq:masters}
  \end{split}
\end{equation}
In these expressions, we used $\csc(z)=1/\sin(z)$, and the
hypergeometric function defined as
\begin{equation}
  \begin{split}
    _2F_1(a,b;c;z) \equiv
    \sum_{n=0}^\infty\frac{(a)_n(b)_n}{(c)_n}\frac{z^n}{n!}\,,
  \end{split}
\end{equation}
with the Pochhammer symbol
\begin{equation}
	(x)_n \equiv \frac{\Gamma ( x + n)}{\Gamma(x)} \,.
\end{equation}
Furthermore, the incomplete beta function is defined by
\begin{equation}
  B_z (a, b) \equiv \int_0^z\dd t\, t^{a-1}(1-t)^{b-1}
\end{equation}
and can be expressed as
\begin{equation}
  \begin{split}
    B_z(a,b) =  z^a \sum_{n = 0}^\infty \frac{(1 - b)_n}{n! (a + n)} z^n
    = \frac{z^a}{a}\ {}_2F_1(a,1-b;a+1;z)\,.
  \end{split}
\end{equation}
The expansions of the hypergeometric function in the limit $\ep\to 0$
can be obtained with the help of the
\texttt{Mathematica}\,\cite{Mathematica11.3} package
\texttt{HypExp}\,\cite{Huber:2005yg,Huber:2007dx}.

A more detailed description of parts of our setup will be described in a
forthcoming publication\,\cite{Artz:2019bpr}. As a check, we evaluated the
correlators $\langle G_{\mu\nu}^aG_{\mu\nu}^a\rangle$, $\langle
\bar\chi\chi\rangle$ and $\langle \bar\chi \slashed D \chi\rangle$
through \nlo. They lead to the same set of master integrals as given in
\eqn{eq:masters}. Comparing our results to \citere{Luscher:2010iy} and
\citere{Makino:2014taa}\footnote{We compare to \texttt{arXiv} versions 2
  and 5 of that paper.}, we find full agreement.
